\documentclass[10pt,a4paper]{amsart}
\usepackage{amsmath,amssymb,amsthm,graphicx}
% Retained minimal preamble for the Stacks Project chapter varieties.tex.
% Verbatim excerpt of the chapter preamble: Xy-pic setup (native preamble lines 14-18), the native theorem
% declarations with the mathematical macros (native preamble lines 154-217), and the project own no-op
% definitions of the reference/slogan/history display environments (native wrapper.tex lines 1-3) that the
% counted statement uses for its slogan. Copied mechanically; nothing added, removed or redefined.
% For commutative diagrams we use Xy-pic
\usepackage[all]{xy}

% We use 2cell for 2-commutative diagrams.
\xyoption{2cell}
% Theorem environments.
%
\theoremstyle{plain}
\newtheorem{theorem}[subsection]{Theorem}
\newtheorem{proposition}[subsection]{Proposition}
\newtheorem{lemma}[subsection]{Lemma}

\theoremstyle{definition}
\newtheorem{definition}[subsection]{Definition}
\newtheorem{example}[subsection]{Example}
\newtheorem{exercise}[subsection]{Exercise}
\newtheorem{situation}[subsection]{Situation}

\theoremstyle{remark}
\newtheorem{remark}[subsection]{Remark}
\newtheorem{remarks}[subsection]{Remarks}

\numberwithin{equation}{subsection}

% Macros
%
\def\lim{\mathop{\mathrm{lim}}\nolimits}
\def\colim{\mathop{\mathrm{colim}}\nolimits}
\def\Spec{\mathop{\mathrm{Spec}}}
\def\Hom{\mathop{\mathrm{Hom}}\nolimits}
\def\Ext{\mathop{\mathrm{Ext}}\nolimits}
\def\SheafHom{\mathop{\mathcal{H}\!\mathit{om}}\nolimits}
\def\SheafExt{\mathop{\mathcal{E}\!\mathit{xt}}\nolimits}
\def\Sch{\mathit{Sch}}
\def\Mor{\mathop{\mathrm{Mor}}\nolimits}
\def\Ob{\mathop{\mathrm{Ob}}\nolimits}
\def\Sh{\mathop{\mathit{Sh}}\nolimits}
\def\NL{\mathop{N\!L}\nolimits}
\def\CH{\mathop{\mathrm{CH}}\nolimits}
\def\proetale{{pro\text{-}\acute{e}tale}}
\def\etale{{\acute{e}tale}}
\def\QCoh{\mathit{QCoh}}
\def\Ker{\mathop{\mathrm{Ker}}}
\def\Im{\mathop{\mathrm{Im}}}
\def\Coker{\mathop{\mathrm{Coker}}}
\def\Coim{\mathop{\mathrm{Coim}}}

% Boxtimes
%
\DeclareMathSymbol{\boxtimes}{\mathbin}{AMSa}{"02}

%
% Macros for moduli stacks/spaces
%
\def\QCohstack{\mathcal{QC}\!\mathit{oh}}
\def\Cohstack{\mathcal{C}\!\mathit{oh}}
\def\Spacesstack{\mathcal{S}\!\mathit{paces}}
\def\Quotfunctor{\mathrm{Quot}}
\def\Hilbfunctor{\mathrm{Hilb}}
\def\Curvesstack{\mathcal{C}\!\mathit{urves}}
\def\Polarizedstack{\mathcal{P}\!\mathit{olarized}}
\def\Complexesstack{\mathcal{C}\!\mathit{omplexes}}
% \Pic is the operator that assigns to X its picard group, usage \Pic(X)
% \Picardstack_{X/B} denotes the Picard stack of X over B
% \Picardfunctor_{X/B} denotes the Picard functor of X over B
\def\Pic{\mathop{\mathrm{Pic}}\nolimits}
\def\Picardstack{\mathcal{P}\!\mathit{ic}}
\def\Picardfunctor{\mathrm{Pic}}
\def\Deformationcategory{\mathcal{D}\!\mathit{ef}}
\newenvironment{reference}{}{}
\newenvironment{slogan}{}{}
\newenvironment{history}{}{}
\usepackage{cleveref}
\crefname{originalreference}{original reference}{original references}
\Crefname{originalreference}{Original reference}{Original references}
\makeatletter
\newlabel{coherent-lemma-flat-base-change-cohomology}{{02KH}{source}{}{sourcepointer-1}{}}
\newlabel{coherent-lemma-flat-base-change-cohomology@cref}{{[originalreference][0][]02KH}{[1][0][]source}{}{}{}}
\makeatother
\crefname{originalreference}{original reference}{original references}
\Crefname{originalreference}{Original reference}{Original references}
\makeatletter
\newlabel{algebra-lemma-Noetherian-irreducible-components}{{00FR}{source}{}{sourcepointer-2}{}}
\newlabel{algebra-lemma-Noetherian-irreducible-components@cref}{{[originalreference][0][]00FR}{[1][0][]source}{}{}{}}
\makeatother
\crefname{originalreference}{original reference}{original references}
\Crefname{originalreference}{Original reference}{Original references}
\makeatletter
\newlabel{schemes-lemma-scheme-finite-discrete-affine}{{02O0}{source}{}{sourcepointer-3}{}}
\newlabel{schemes-lemma-scheme-finite-discrete-affine@cref}{{[originalreference][0][]02O0}{[1][0][]source}{}{}{}}
\makeatother
\crefname{originalreference}{original reference}{original references}
\Crefname{originalreference}{Original reference}{Original references}
\makeatletter
\newlabel{coherent-lemma-relative-affine-vanishing}{{01XC}{source}{}{sourcepointer-4}{}}
\newlabel{coherent-lemma-relative-affine-vanishing@cref}{{[originalreference][0][]01XC}{[1][0][]source}{}{}{}}
\makeatother
\title{Vanishing of the higher direct images along the inclusion of a quasi-compact open in a one-dimensional affine Noetherian case}
\author{The Stacks Project Authors}
\date{stacks-project:varieties.tex:local-snapshot-20261011; GNU FDL 1.2 or later}
\begin{document}
\maketitle
\noindent\textit{Editorial note.} All delivered mathematical text is the author's own, copied verbatim from the pinned Stacks Project chapter file varieties.tex; no statement or proof step was rewritten, reordered or removed.\par\medskip
\noindent\textit{Editorial note.} The counted span cites four cross-chapter labels, each recorded in assembly.external\_references with its canonical tag as a navigation pointer.\par\medskip
\noindent\textit{Editorial note.} The proof cites no bibliography item inside the counted span.\par\medskip
\noindent\textit{Allowed standard prerequisites (editorial).}\begin{itemize}
\item Commutation of higher direct images with flat base change. Stacks Project, Cohomology of Schemes, Lemma coherent-lemma-flat-base-change-cohomology (tag 02KH). Justifies reducing the stalkwise vanishing to the case of a local ring.
\item A finite discrete scheme is affine; irreducible components of a Noetherian local ring. Stacks Project, Schemes, Lemma schemes-lemma-scheme-finite-discrete-affine (tag 0CKF); Algebra, Lemma algebra-lemma-Noetherian-irreducible-components (tag 00FZ). Supply the explicit description of U and its affineness in the local case.
\end{itemize}\medskip
\ifdefined\hypertarget\hypertarget{sourcepointer-1}{}\fi
\noindent\textit{Source reference.} Original reference 02KH: sources/stacks-chow/tags.txt, source lines 3372--3372. Cross-chapter standard prerequisite cited by the author inside the counted span; it is used as a pointer and its own proof is not part of this item.\par\medskip
\ifdefined\hypertarget\hypertarget{sourcepointer-2}{}\fi
\noindent\textit{Source reference.} Original reference 00FR: sources/stacks-chow/tags.txt, source lines 581--581. Cross-chapter standard prerequisite cited by the author inside the counted span; it is used as a pointer and its own proof is not part of this item.\par\medskip
\ifdefined\hypertarget\hypertarget{sourcepointer-3}{}\fi
\noindent\textit{Source reference.} Original reference 02O0: sources/stacks-chow/tags.txt, source lines 3499--3499. Cross-chapter standard prerequisite cited by the author inside the counted span; it is used as a pointer and its own proof is not part of this item.\par\medskip
\ifdefined\hypertarget\hypertarget{sourcepointer-4}{}\fi
\noindent\textit{Source reference.} Original reference 01XC: sources/stacks-chow/tags.txt, source lines 2513--2513. Cross-chapter standard prerequisite cited by the author inside the counted span; it is used as a pointer and its own proof is not part of this item.\par\medskip
\begin{lemma}
\label{lemma-affine}
Let $X$ be a scheme all of whose local rings are Noetherian of dimension
$\leq 1$. Let $U \subset X$ be a retrocompact open. Denote
$j : U \to X$ the inclusion morphism. Then $R^pj_*\mathcal{F} = 0$, $p > 0$
for every quasi-coherent $\mathcal{O}_U$-module $\mathcal{F}$.
\end{lemma}

\begin{proof}
We may check the vanishing of $R^pj_*\mathcal{F}$ at stalks.
Formation of $R^qj_*$ commutes with flat base change, see
Cohomology of Schemes, Lemma
\ref{coherent-lemma-flat-base-change-cohomology}.
Thus we may assume that $X$ is the spectrum of a Noetherian local
ring of dimension $\leq 1$. In this case $X$ has a closed point
$x$ and finitely many other points $x_1, \ldots, x_n$ which specialize
to $x$ but not each other (see
Algebra, Lemma \ref{algebra-lemma-Noetherian-irreducible-components}).
If $x \in U$, then $U = X$ and the result is clear. If not, then
$U = \{x_1, \ldots, x_r\}$ for some $r$ after possibly renumbering
the points. Then $U$ is affine
(Schemes, Lemma \ref{schemes-lemma-scheme-finite-discrete-affine}).
Thus the result follows from Cohomology of Schemes, Lemma
\ref{coherent-lemma-relative-affine-vanishing}.
\end{proof}

\end{document}
